%% kapa: exam=1 round=33 year=2022 period=1 form=A subject=economics subject_ko=경제학원론 seq=9 number=49
%% source: Q-Net 감정평가사 기출문제 (cst00309) artlSeq=5209997 file=raw/1008/5209997/2022년도 제33회 감정평가사 1차시험 1교시(민법,경제학,부동산학).pdf
%% conversion: kordoc 4.24.51 -> tools/gen1.py (auto); review=proofread-sonnet; flags=none
\begin{question}{49}
( )에 들어갈 내용으로 옳은 것은?
\begin{framedtext}
과점시장에서 보수를 극대화하는 두 기업 A와 B가 각각 전략 1과 전략 2를 통해 아래 표와 같은 보수(payoff)를 얻을 수 있다. \par
{\small\centering\begin{tabular}{cc|c|c|}
\multicolumn{2}{c}{} & \multicolumn{2}{c}{기업 B} \\ \cline{3-4}
\multicolumn{2}{c|}{} & 전략 1 & 전략 2 \\ \cline{2-4}
\multirow{2}{*}{기업 A} & 전략 1 & (22, 10) & (33, 8) \\ \cline{2-4}
 & 전략 2 & (32, 14) & (30, 12) \\ \cline{2-4}
\end{tabular}\par}
\hfill ※ ( )안의 앞의 숫자는 기업 A의 보수, 뒤의 숫자는 기업 B의 보수이다. \par
\smallskip
○ 기업 A와 기업 B가 동시에 전략을 선택할 때, 균형에서 기업 A의 보수는 ( ㄱ )이다. \par
○ 기업 A가 먼저 전략을 선택하고 신뢰할 수 있는 방법으로 확약할 때, 균형에서 기업 B의 보수는 ( ㄴ )이다.
\end{framedtext}
\begin{choices}
  \item ㄱ: 22, ㄴ: 8
  \item ㄱ: 30, ㄴ: 8
  \item ㄱ: 32, ㄴ: 10
  \item ㄱ: 32, ㄴ: 14
  \item ㄱ: 33, ㄴ: 12
\end{choices}
\end{question}
\answer{4}
